\documentclass[11pt]{article}
\usepackage[margin=1in]{geometry}
\usepackage{amsmath,amssymb}
\usepackage{booktabs}
\usepackage{array}
\usepackage{xcolor}
\usepackage[colorlinks=true,linkcolor=blue,urlcolor=blue]{hyperref}
\usepackage{enumitem}

\newcommand{\code}[1]{\texttt{\small #1}}
\newcommand{\env}[1]{\texttt{\small #1}}

\title{\textbf{SPAX RVE Toolkit --- New Features}\\[2pt]
\large Quadratic couple-stress homogenization, anisotropy reporting,\\
and tool consolidation}
\author{Spatium development notes}
\date{June 2026}

\begin{document}
\maketitle

\begin{abstract}
This note documents the features added to the SPAX representative-volume-element
(RVE) homogenization toolkit. The headline additions are: a second-order
(quadratic, \code{C3D10H}) periodic bending pipeline that removes the linear-tet
shear-locking that previously contaminated the apparent couple-stress response;
a per-phase element formulation that keeps hybrid elements only where they are
needed; corrected volume-weighted post-processing for multi-integration-point
elements; full effective-anisotropy reporting including the $6{\times}6$
elasticity tensor; a permanent implementation of the Choi et al. (2022) Eq.~(19)
modified-couple-stress (MCST) length-scale extraction; and consolidation of the
entire workflow into three Python files.
\end{abstract}

\tableofcontents

%==============================================================
\section{Overview}
The toolkit is now contained in three files:
\begin{itemize}[nosep]
  \item \code{Spatium\_Standalone.py} --- packing, periodic meshing driver,
        Abaqus input-deck (\code{.inp}) generation;
  \item \code{Spatium\_GmshPeriodic.py} --- the periodic Gmsh mesher;
  \item \code{Spatium\_PostProcess.py} --- ODB extraction, effective properties,
        anisotropy, and all study-level analyzers.
\end{itemize}
All new behaviour is opt-in through environment variables (Table~\ref{tab:env});
defaults reproduce the previous (linear, isotropic-report) pipeline exactly.

%==============================================================
\section{Second-order (quadratic) bending elements}
\label{sec:quad}

\subsection{Motivation}
Linear tetrahedra (\code{C3D4H}) reproduce a uniform strain field exactly but
shear-lock severely in bending. On a homogeneous cube (which has, by
construction, \emph{no} intrinsic length scale) the apparent bending stiffness
ratio
\begin{equation}
  f(L)=\frac{D_\text{rve}}{E_\text{plate}\,L^4/12},
  \qquad E_\text{plate}=\frac{E}{1-\nu^2},
\end{equation}
should equal a size-independent constant. With linear elements $f$ ranged from
$1.12$ to $1.28$ (worse for coarse, bending-starved meshes); with quadratic
\code{C3D10H} elements $f\!\approx\!1.0$--$1.08$. The spurious $+12$ to $+28\%$
was therefore element shear-locking, not physics.

\subsection{The \code{SPAX\_MESH\_ORDER=2} pipeline}
Setting \env{SPAX\_MESH\_ORDER=2} activates a quadratic mesh. Because the custom
periodic mesher enforces conformity by copying master face meshes onto slaves,
order promotion is performed \emph{after} the linear periodic volume mesh is
complete:
\begin{enumerate}[nosep]
  \item generate the conforming linear periodic tetrahedral mesh as usual;
  \item \code{gmsh.option Mesh.SecondOrderLinear=1} --- keep elements
        \emph{straight-sided} (mid-edge nodes at edge midpoints). This avoids
        inverting coarse elements on curved inclusion surfaces and keeps the
        outer-face mid-edge nodes exact periodic translates;
  \item \code{gmsh.model.mesh.setOrder(2)};
  \item a \emph{validity gate}: the minimum scaled-inverse-condition-number
        (\code{minSICN}) of all 3-D elements is checked; any inverted/degenerate
        element raises an exception that triggers a re-pack of the
        microstructure (slivers are stochastic), repeating until a clean mesh is
        obtained (\env{SPAX\_MIN\_SICN}, default $0.01$).
\end{enumerate}
The coordinate-based periodic-pair matcher then pairs the new mid-edge boundary
nodes automatically (zero skipped pairs).

\subsection{Gmsh\,$\rightarrow$\,Abaqus C3D10 node remap}
Gmsh orders the last two mid-edge nodes of a ten-node tetrahedron opposite to
the Abaqus \code{C3D10} convention. Without correction every element is read as
zero/negative volume and the input processor aborts. The deck writer swaps the
9th and 10th connectivity entries,
\[
\underbrace{[c_0,c_1,c_2,c_3,\,m_{01},m_{12},m_{20},m_{03},\,m_{23},m_{13}]}_{\text{Gmsh}}
\;\longrightarrow\;
\underbrace{[c_0,c_1,c_2,c_3,\,m_{01},m_{12},m_{20},m_{03},\,m_{13},m_{23}]}_{\text{Abaqus C3D10}} .
\]

\subsection{Cap rejection for clean quadratic meshes}
For order-2 runs the packer rejects inclusion placements that only graze a
periodic face (thin caps that mesh as positive-volume linear tets but invert as
quadratic ones). A constant cap-rejection floor of
$\env{SPAX\_SLIVER\_MULT\_Q}\times\code{lc\_fine}$ (default $1.0$) is enforced
from the first attempt, eliminating the face slivers at the default element size
($\sim 16$\,k elements for an $L/d{=}3$ box, versus $\sim 40$\,k from brute-force
global refinement).

%==============================================================
\section{Per-phase (heterogeneous) element formulation}
\label{sec:hybrid}
The hybrid (\code{H}) formulation adds an internal pressure degree of freedom
that prevents volumetric locking of (near-)incompressible material, at the cost
of a larger linear system. Only the brine inclusions ($\nu\!\approx\!0.49$)
require it; the compressible ice matrix ($\nu\!\approx\!0.33$) does not, and the
matrix is the majority of elements. The deck writer now selects the element type
\emph{per phase}, keyed on Poisson's ratio:
\begin{equation}
  \text{type}(\nu)=
  \begin{cases}
     \code{C3D10H}/\code{C3D4H} & \nu \ge \nu_{H}\ (\text{inclusions}),\\[2pt]
     \code{C3D10}/\code{C3D4}   & \nu < \nu_{H}\ (\text{matrix}),
  \end{cases}
  \qquad \nu_H=\env{SPAX\_HYBRID\_NU}\ (\text{default }0.45).
\end{equation}
On a representative $L/d{=}4$ porous bending solve this reduced the system from
$247{,}549$ to $222{,}640$ variables ($-10\%$) and the wall-clock solve time by
$13\%$, changing $D_\text{rve}$ by $<\!1\%$ (which cancels in the
homogeneous-calibrated ratio). \env{SPAX\_FORCE\_HYBRID=1} restores the legacy
all-hybrid behaviour.

%==============================================================
\section{Post-processing correction for multi-IP elements}
\label{sec:multiip}
A quadratic tetrahedron carries four integration points (IPs); a linear one
carries a single whole-element value. The volume-weighted averages
$\langle\sigma\rangle = \sum_e \sigma_e V_e / V_\text{RVE}$ and the bending
moment were written for one value per element. Read na\"ively for \code{C3D10},
each of an element's four IP rows was weighted by the \emph{full} element volume
\code{EVOL}, a four-fold over-count of $D_\text{rve}$, $E_\text{eff}$ and the
solid volume (porosity even went negative). The fix collapses IP rows to their
per-element mean before weighting,
\begin{equation}
  \bar{\sigma}_e=\frac{1}{n_\text{IP}}\sum_{p=1}^{n_\text{IP}}\sigma_{e,p},
\end{equation}
inside \code{\_field\_by\_label} (\code{\_mean\_per\_element}); it is a no-op for
single-IP \code{C3D4}. Note that \emph{dimensionless ratios} ($f$,
$E_\text{eff}=\sigma/\varepsilon$, $D/D_\text{class}$) are unaffected by the bug
because the factor cancels; only absolute cross-run comparisons required the fix.

%==============================================================
\section{Effective anisotropy and the elasticity tensor}
\label{sec:aniso}
For a full-tensor run (\env{full\_tensor=Yes}, six uniaxial + shear load cases)
the post-processor now reports the \emph{directional} effective moduli rather
than a single scalar:
\[
  E_x,E_y,E_z,\quad G_{xy},G_{xz},G_{yz},\quad \nu_x,\nu_y,\nu_z,
\]
together with anisotropy indices
\begin{equation}
  A=\frac{\max(E_x,E_y,E_z)}{\min(E_x,E_y,E_z)},
  \qquad
  R_{z/xy}=\frac{E_z}{\tfrac12(E_x+E_y)} .
\end{equation}
$R_{z/xy}$ targets the transverse isotropy of the vertical ($z$) brine-channel
network. A representative channel RVE returned
$E_x{=}6.12$, $E_y{=}6.07$, $E_z{=}7.68$~GPa, i.e.\ $R_{z/xy}=1.26$ ---
$26\%$ stiffer along the channel axis, the columnar-ice signature. The scalar
$E_\text{eff}/\nu_\text{eff}/G_\text{eff}$ are retained (taken from the declared
\code{Mode}/\code{Mode2}) for backward compatibility.

The complete anisotropy descriptor, the $6\times6$ effective stiffness
$C_{ij}$, is produced by the \code{elasticity} sub-command, which assembles the
columns of $C$ from the volume-averaged stress response to the six unit
macro-strains.

%==============================================================
\section{MCST length-scale extraction (Eq.~19)}
\label{sec:eq19}
The modified-couple-stress size effect is, following Choi, Lee \& Sim
(\textit{Mater.\ Des.}\ \textbf{214}, 2022, Eq.~19),
\begin{equation}
  E_\text{app}(L)=E_\text{classical}+12\,\mu\left(\frac{\ell}{L}\right)^{2},
  \qquad E_\text{app}=\frac{12\,D_\text{rve}}{L^{4}} .
  \label{eq:eq19}
\end{equation}
Fitting $E_\text{app}$ against $1/L^2$ over a size sweep yields the length scale
from the \emph{slope},
\begin{equation}
  \ell=\sqrt{\frac{\text{slope}}{12\,\mu}},
\end{equation}
which is \emph{reference-free}: the intercept absorbs the (plate vs.\ beam)
classical modulus, so $\ell$ does not depend on that choice. The intercept itself
is a useful consistency check --- for our RVEs it lands on $E^\ast/(1-\nu^2)$,
confirming plate-like bending.

The analyzer also reports, per realization, the bending-vs-first-order modulus
ratio (geometric plate factor removed),
\begin{equation}
  \frac{E_\text{bend}^\text{mat}}{E^\ast_\text{1st}},
  \qquad E_\text{bend}^\text{mat}=E_\text{app}\,(1-\nu^2),
\end{equation}
separating first-order softening (from weak inclusions) from any genuine MCST
stiffening. A ratio $\approx 1$ with no size trend indicates no length scale;
a \emph{size-independent} excess indicates classical anisotropy rather than
couple stress. Applied to the sea-ice microstructures (isolated soft pockets and
vertical brine channels) the slope is non-positive in both cases: \emph{no
measurable couple-stress length scale}; the channel excess is the classical
transverse isotropy of Sec.~\ref{sec:aniso}.

%==============================================================
\section{Consolidated tool and command-line interface}
\label{sec:cli}
All former standalone analysis scripts are folded into
\code{Spatium\_PostProcess.py}. \code{odbAccess} is imported lazily, so the
CSV-only analyzers run under a plain \code{python3} with no Abaqus licence.

\begin{center}
\begin{tabular}{@{}>{\raggedright\arraybackslash}p{0.46\linewidth} >{\raggedright\arraybackslash}p{0.46\linewidth}@{}}
\toprule
\textbf{Command} & \textbf{Purpose}\\
\midrule
\code{abaqus python \dots\ <csv> <odb\_dir> [out] [idx]} & per-RVE extraction (incl.\ anisotropy)\\
\code{python \dots\ --merge <parts> <out>} & union per-RVE partial CSVs\\
\code{python \dots\ analyze eq19 <csv\ldots>} & Eq.~19 MCST fit + bend/1st-order\\
\code{python \dots\ analyze lengthscale <por> [base\ldots]} & homogeneous-calibrated $\ell$\\
\code{python \dots\ analyze homog-calib} & homogeneous-cube artifact $f(L)$\\
\code{python \dots\ analyze hybrid <fo> <bend>} & small-RVE moduli $+$ big-RVE $D_\text{rve}$\\
\code{python \dots\ analyze rve-study <csv>} & RVE-size convergence (mean/CoV)\\
\code{abaqus python \dots\ principal <odb> <out>} & principal stresses/strains/energy\\
\code{abaqus python \dots\ elasticity <dir> <out> <L> <id>} & full $6\times6$ $C_{ij}$\\
\bottomrule
\end{tabular}
\end{center}

%==============================================================
\section{Environment-variable reference}
\begin{table}[h]
\centering
\small
\begin{tabular}{@{}lll@{}}
\toprule
\textbf{Variable} & \textbf{Default} & \textbf{Effect}\\
\midrule
\env{SPAX\_MESH\_ORDER}     & \code{1} & \code{2} $\Rightarrow$ quadratic \code{C3D10(H)}\\
\env{SPAX\_MIN\_SICN}       & \code{0.01} & quadratic element validity floor (re-pack gate)\\
\env{SPAX\_SLIVER\_MULT\_Q} & \code{1.0} & order-2 cap-rejection floor ($\times$\,\code{lc\_fine})\\
\env{SPAX\_OPT\_PASSES}     & \code{2/3} & Netgen quality passes before promotion\\
\env{SPAX\_HYBRID\_NU}      & \code{0.45} & $\nu$ threshold for hybrid elements\\
\env{SPAX\_FORCE\_HYBRID}   & \code{0} & \code{1} $\Rightarrow$ legacy all-hybrid\\
\env{SPAX\_MESH\_ALGO2D}    & (Gmsh) & optional 2-D mesher override\\
\env{SPAX\_SEED}            & (entropy) & reproducible packing (deterministic per row)\\
\env{SPAX\_SAVE\_PACKING}   & --- & freeze packed sphere array to \code{<dir>/<run\_id>.npy}\\
\env{SPAX\_LOAD\_PACKING}   & --- & reuse a frozen packing (skip packer); re-mesh same geometry\\
\bottomrule
\end{tabular}
\end{table}

%==============================================================
\section{Validation summary}
\begin{itemize}[nosep]
  \item Homogeneous cube ($\ell\!=\!0$ by construction): quadratic $f\!\to\!1.0$
        (vs.\ $1.12$--$1.28$ linear), confirming the locking removal.
  \item Per-phase elements: $D_\text{rve}$ within $1\%$ of all-hybrid at
        $-13\%$ solve time.
  \item Multi-IP fix: homogeneous $f$ unchanged (it is a ratio); porous absolute
        $D_\text{rve}$, $E_\text{eff}$ and porosity restored to physical values.
  \item Mesh resolution: with the geometry \emph{frozen}
        (\env{SPAX\_LOAD\_PACKING}), porous $D_\text{rve}$ is converged --- it
        varies only $\sim\!1\%$ across $\code{L\_mesh}=0.045\to0.026$. An earlier
        ``same seed'' study had wrongly suggested $+30\%$ sensitivity; that was a
        \emph{packing} confound (the packer's $r_\text{floor}$ and sliver floors
        scale with \code{L\_mesh}, so changing the mesh also changed the
        inclusion arrangement at fixed porosity). Only freezing the actual sphere
        array isolates the mesh effect, which is small.
  \item Eq.~19 reproduces the standalone analyzer exactly and, across all
        methods tried, returns no positive couple-stress slope for either
        microstructure.
\end{itemize}

\end{document}
